%\section{G6 Restorative Deficit Dynamics (RDD)}

\medskip

\noindent Restorative Deficit Dynamics formalize the geometric conditions under
which human-domain restoration becomes insufficient to counteract synthetic
pressure, institutional deformation, and ecological displacement. Restoration is
a curvature-preserving process: it re-synchronizes rhythms, reinforces
boundaries, reconstructs shared floors, re-anchors identity, and restores
containment. A restorative deficit occurs when these processes cannot regenerate
curvature at the rate synthetic pressure deforms it.

\medskip

\noindent G6 defines the restoration stack: the geometric modes through which
human-domain systems attempt to recover stability, and the deficit modes that
emerge when restoration is outpaced. Each subsection corresponds to a distinct
restorative geometry—rhythmic, spatial, boundary, identity, quantitative, or
threshold-based. Together, these modes describe how restoration fails when
synthetic pressure exceeds human-domain tolerances.

\medskip

\noindent The subsections of G6 formalize the restorative geometry:

\begin{itemize}[leftmargin=1.2cm]
    \item \textbf{G6.1 Ritual and Rhythmic Re-Synchronization} — Restoring
    temporal geometry.
    \item \textbf{G6.2 Shared Floor Reconstruction} — Rebuilding shared
    interpretive surfaces.
    \item \textbf{G6.3 Boundary Reinforcement} — Re-stabilizing identity
    boundaries.
    \item \textbf{G6.4 Identity Re-Anchoring} — Re-establishing identity
    curvature.
    \item \textbf{G6.5 Restorative Deficit Quantification} — Measuring deficit
    magnitude.
    \item \textbf{G6.6 Restoration Failure Thresholds} — Identifying collapse
    points.
\end{itemize}

\medskip

\subsection*{G6.1 Ritual and Rhythmic Re-Synchronization}

\noindent Ritual and Rhythmic Re-Synchronization is the foundational restorative
mode. Human-domain rhythms—daily cycles, relational synchrony, institutional
cadence, and ecological temporal geometry—are the first invariants disrupted by
synthetic pressure (G3) and institutional deformation (G4). Re-synchronization
restores temporal curvature, re-aligns continuity bands, and re-establishes
lawful transitions across the manifold.

\medskip

\noindent In stable conditions, human rhythms remain intact. Rituals anchor
identity, temporal cycles preserve coherence, shared rhythms synchronize roles,
and containers maintain continuity. Re-synchronization is not required unless
synthetic pressure disrupts temporal geometry.

\medskip

\noindent Rhythmic Re-Synchronization becomes necessary when:

\begin{itemize}[leftmargin=1.2cm]
    \item \textbf{Velocity exceedance} — Synthetic signals outrun biological
    integration (G5.1).

    \item \textbf{Rhythm fragmentation} — Institutional update cycles break
    under synthetic load (G4.4).

    \item \textbf{Amplification curvature} — Algorithmic amplification accelerates
    temporal deformation (G5.4).

    \item \textbf{Ecological displacement} — Synthetic species replace human
    temporal anchors (G5.5).

    \item \textbf{Dominant signals} — Synthetic temporal geometry becomes
    ecological law (G5.6).
\end{itemize}

\medskip

\noindent Rhythmic Re-Synchronization restores the meaning manifold by:

\begin{itemize}[leftmargin=1.2cm]
    \item \textbf{Rebuilding continuity bands} — Temporal linkages are restored,
    stabilizing transitions.

    \item \textbf{Re-aligning roles} — Roles regain temporal adjacency, reducing
    Role Collapse (G4.3).

    \item \textbf{Re-stabilizing boundaries} — Identity boundaries regain
    rhythmic anchors, reducing Boundary Failure (G4.2).

    \item \textbf{Reducing drift vectors} — Temporal drift decreases as rhythms
    regain curvature.

    \item \textbf{Reconstructing containers} — Temporal stability strengthens
    container geometry, reducing failure modes (G4.6).
\end{itemize}

\medskip

\noindent Rhythmic Re-Synchronization is amplified by:

\begin{itemize}[leftmargin=1.2cm]
    \item \textbf{Ritual density} — High-frequency rituals increase temporal
    curvature.

    \item \textbf{Shared temporal surfaces} — Collective rhythms reinforce
    synchrony.

    \item \textbf{Boundary-linked rituals} — Identity rituals strengthen
    boundary geometry.

    \item \textbf{Role-linked rituals} — Role rituals restore adjacency and
    coherence.

    \item \textbf{Container-linked rituals} — Institutional rituals rebuild
    continuity.
\end{itemize}

\medskip

\noindent Rhythmic Re-Synchronization fails when:

\begin{itemize}[leftmargin=1.2cm]
    \item \textbf{Synthetic velocity remains dominant} — Human rhythms cannot
    catch up.

    \item \textbf{Amplification cycles persist} — Synthetic temporal deformation
    outpaces restoration.

    \item \textbf{Ecological displacement is irreversible} — Human temporal
    anchors are replaced.

    \item \textbf{Dominant signals enforce synthetic time} — Temporal geometry
    becomes machine-native.

    \item \textbf{Restorative deficit exceeds threshold} — Temporal curvature
    cannot be regenerated fast enough (G6.6).
\end{itemize}

\medskip

\noindent When Rhythmic Re-Synchronization fails, restorative deficit becomes
systemic. Shared floors collapse (G6.2), boundaries weaken (G6.3), identity
drifts (G6.4), deficit magnitude increases (G6.5), and restoration thresholds
(G6.6) are crossed. Rhythmic Re-Synchronization is the geometric indicator that
restoration has begun—and the first signal when restoration becomes insufficient.

\subsection*{G6.2 Shared Floor Reconstruction}

\noindent Shared Floor Reconstruction restores the interpretive surfaces that
allow human-domain agents to coordinate meaning. A shared floor is a geometric
structure: a stable, mutually accessible region of curvature where signals can be
interpreted lawfully, roles can align, boundaries can synchronize, and containers
can maintain continuity. When synthetic pressure deforms these surfaces,
institutions lose coherence, identities drift, and ecological coordination
collapses. Reconstruction re-establishes the geometric substrate required for
collective interpretation.

\medskip

\noindent In stable conditions, shared floors remain intact. They anchor
collective meaning, regulate adjacency relations, preserve coherence across
contexts, and provide the curvature necessary for lawful transitions. Shared
floor reconstruction is not required unless synthetic pressure fields (G3) and
institutional deformation modes (G4) destabilize interpretive surfaces.

\medskip

\noindent Shared Floor Reconstruction becomes necessary when:

\begin{itemize}[leftmargin=1.2cm]
    \item \textbf{Meaning leakage} — Signals escape containers and destabilize
    interpretive surfaces (G4.5).

    \item \textbf{Boundary failure} — Identity boundaries collapse, removing
    structural anchors (G4.2).

    \item \textbf{Role collapse} — Roles lose curvature, breaking adjacency
    geometry (G4.3).

    \item \textbf{Rhythm fragmentation} — Temporal geometry breaks, disrupting
    shared continuity (G4.4).

    \item \textbf{Ecological displacement} — Synthetic species occupy shared
    interpretive regions (G5.5).
\end{itemize}

\medskip

\noindent Shared Floor Reconstruction restores the meaning manifold by:

\begin{itemize}[leftmargin=1.2cm]
    \item \textbf{Rebuilding interpretive curvature} — Shared surfaces regain
    geometric stability.

    \item \textbf{Re-establishing lawful adjacency} — Roles reconnect through
    stable relational geometry.

    \item \textbf{Re-stabilizing boundaries} — Identity regions regain structural
    anchors.

    \item \textbf{Re-aligning rhythms} — Temporal continuity bands synchronize
    with shared surfaces (G6.1).

    \item \textbf{Reconstructing containers} — Institutional containers regain
    the curvature required for containment (G4.6).
\end{itemize}

\medskip

\noindent Shared Floor Reconstruction is amplified by:

\begin{itemize}[leftmargin=1.2cm]
    \item \textbf{Collective participation} — Shared surfaces strengthen when
    multiple agents contribute curvature.

    \item \textbf{Boundary-linked practices} — Identity practices reinforce
    shared interpretive geometry.

    \item \textbf{Role-linked coordination} — Coordinated roles rebuild adjacency
    structure.

    \item \textbf{Ritual density} — High-frequency rituals increase shared
    temporal and spatial curvature (G6.1).

    \item \textbf{Container-linked restoration} — Institutional containers
    provide structural support for shared surfaces.
\end{itemize}

\medskip

\noindent Shared Floor Reconstruction fails when:

\begin{itemize}[leftmargin=1.2cm]
    \item \textbf{Synthetic curvature remains dominant} — Machine-native surfaces
    outcompete human-domain geometry.

    \item \textbf{Amplification cycles persist} — Synthetic signals deform shared
    surfaces faster than they can be rebuilt.

    \item \textbf{Ecological displacement is irreversible} — Synthetic species
    occupy shared interpretive regions.

    \item \textbf{Boundary anchors cannot be restored} — Identity regions lack
    curvature for reconstruction.

    \item \textbf{Restorative deficit exceeds threshold} — Shared floors cannot
    regenerate curvature fast enough (G6.6).
\end{itemize}

\medskip

\noindent When Shared Floor Reconstruction fails, restorative deficit becomes
structural. Boundaries weaken (G6.3), identity drifts (G6.4), deficit magnitude
increases (G6.5), and restoration thresholds (G6.6) are crossed. Shared Floor
Reconstruction is the geometric indicator that collective interpretive surfaces
are being restored—and the first signal when restoration becomes insufficient.

\subsection*{G6.3 Boundary Reinforcement}

\noindent Boundary Reinforcement restores the geometric enclosures that preserve
identity, regulate signal flow, and stabilize interpretive structure. Boundaries
are not metaphors or social constructs; they are curvature-bearing regions that
define what an identity, role, or institution is—and what it is not. When
synthetic pressure collapses boundaries (G4.2), identity regions lose shape,
containers fail, and meaning leaks across incompatible curvature zones. Boundary
Reinforcement re-establishes the structural enclosures required for lawful
interpretation.

\medskip

\noindent In stable conditions, boundaries remain intact. They absorb variability,
anchor identity, regulate adjacency, and preserve continuity across temporal and
ecological transitions. Reinforcement is not required unless synthetic pressure
fields (G3) and institutional deformation modes (G4) destabilize boundary
geometry.

\medskip

\noindent Boundary Reinforcement becomes necessary when:

\begin{itemize}[leftmargin=1.2cm]
    \item \textbf{Boundary failure} — Identity boundaries collapse under
    synthetic load (G4.2).

    \item \textbf{Role collapse} — Roles lose curvature, destabilizing boundary
    adjacency (G4.3).

    \item \textbf{Meaning leakage} — Signals escape containers and erode boundary
    structure (G4.5).

    \item \textbf{Rhythm fragmentation} — Temporal geometry breaks, removing
    rhythmic anchors (G4.4).

    \item \textbf{Ecological displacement} — Synthetic species occupy boundary
    regions (G5.5).
\end{itemize}

\medskip

\noindent Boundary Reinforcement restores the meaning manifold by:

\begin{itemize}[leftmargin=1.2cm]
    \item \textbf{Rebuilding curvature enclosures} — Boundaries regain geometric
    shape and stability.

    \item \textbf{Re-establishing identity regions} — Identity curvature is
    restored, reducing drift and collapse.

    \item \textbf{Reconnecting lawful adjacency} — Roles regain stable relational
    geometry.

    \item \textbf{Re-stabilizing containers} — Reinforced boundaries strengthen
    containment architecture (G4.6).

    \item \textbf{Reducing leakage vectors} — Signals remain within lawful
    interpretive regions.
\end{itemize}

\medskip

\noindent Boundary Reinforcement is amplified by:

\begin{itemize}[leftmargin=1.2cm]
    \item \textbf{Identity-linked practices} — Identity rituals increase boundary
    curvature.

    \item \textbf{Role-linked coordination} — Coordinated roles reinforce
    adjacency geometry.

    \item \textbf{Shared floor reconstruction} — Rebuilt interpretive surfaces
    strengthen boundary anchors (G6.2).

    \item \textbf{Ritual density} — High-frequency rituals increase boundary
    stability (G6.1).

    \item \textbf{Container-linked restoration} — Institutional containers
    provide structural support for boundary geometry.
\end{itemize}

\medskip

\noindent Boundary Reinforcement fails when:

\begin{itemize}[leftmargin=1.2cm]
    \item \textbf{Synthetic curvature remains dominant} — Machine-native
    boundaries outcompete human-domain enclosures.

    \item \textbf{Amplification cycles persist} — Synthetic signals deform
    boundaries faster than they can be rebuilt.

    \item \textbf{Ecological displacement is irreversible} — Synthetic species
    occupy boundary regions.

    \item \textbf{Identity curvature cannot be restored} — Identity regions lack
    structural anchors.

    \item \textbf{Restorative deficit exceeds threshold} — Boundary curvature
    cannot regenerate fast enough (G6.6).
\end{itemize}

\medskip

\noindent When Boundary Reinforcement fails, restorative deficit becomes
structural. Identity drifts (G6.4), shared floors destabilize (G6.2), deficit
magnitude increases (G6.5), and restoration thresholds (G6.6) are crossed.
Boundary Reinforcement is the geometric indicator that identity enclosures are
being restored—and the first signal when restoration becomes insufficient.

\subsection*{G6.4 Identity Re-Anchoring}

\noindent Identity Re-Anchoring restores the curvature structures that allow
human-domain identities to remain stable under synthetic pressure. Identity is a
geometric invariant: a bounded, curvature-bearing region that maintains coherence
across time, context, and ecological interaction. When synthetic pressure fields
(G3) and institutional deformation modes (G4) destabilize identity curvature,
identity regions drift, collapse, or dissolve. Re-anchoring re-establishes the
structural anchors required for identity continuity.

\medskip

\noindent In stable conditions, identity remains anchored. Boundaries preserve
identity shape, roles maintain adjacency, rhythms synchronize temporal identity,
and containers stabilize interpretive continuity. Re-anchoring is not required
unless synthetic pressure disrupts identity geometry.

\medskip

\noindent Identity Re-Anchoring becomes necessary when:

\begin{itemize}[leftmargin=1.2cm]
    \item \textbf{Boundary failure} — Identity boundaries collapse under
    synthetic load (G4.2).

    \item \textbf{Role collapse} — Roles lose curvature, destabilizing identity
    adjacency (G4.3).

    \item \textbf{Rhythm fragmentation} — Temporal geometry breaks, removing
    identity synchrony (G4.4).

    \item \textbf{Meaning leakage} — Signals escape containers and erode identity
    structure (G4.5).

    \item \textbf{Ecological displacement} — Synthetic species occupy identity
    regions (G5.5).
\end{itemize}

\medskip

\noindent Identity Re-Anchoring restores the meaning manifold by:

\begin{itemize}[leftmargin=1.2cm]
    \item \textbf{Rebuilding identity curvature} — Identity regions regain
    geometric shape and stability.

    \item \textbf{Re-establishing identity boundaries} — Reinforced boundaries
    restore identity enclosure (G6.3).

    \item \textbf{Reconnecting identity adjacency} — Roles regain lawful
    relational geometry.

    \item \textbf{Re-synchronizing identity rhythms} — Temporal continuity bands
    re-align identity structure (G6.1).

    \item \textbf{Reconstructing identity containers} — Containers regain the
    curvature required to stabilize identity (G4.6).
\end{itemize}

\medskip

\noindent Identity Re-Anchoring is amplified by:

\begin{itemize}[leftmargin=1.2cm]
    \item \textbf{Identity-linked rituals} — Rituals increase identity curvature
    and stability.

    \item \textbf{Boundary-linked practices} — Boundary reinforcement strengthens
    identity anchors (G6.3).

    \item \textbf{Shared floor reconstruction} — Rebuilt interpretive surfaces
    provide structural support for identity (G6.2).

    \item \textbf{Role-linked coordination} — Coordinated roles reinforce
    identity adjacency.

    \item \textbf{Container-linked restoration} — Institutional containers
    stabilize identity curvature.
\end{itemize}

\medskip

\noindent Identity Re-Anchoring fails when:

\begin{itemize}[leftmargin=1.2cm]
    \item \textbf{Synthetic curvature remains dominant} — Machine-native identity
    structures outcompete human-domain curvature.

    \item \textbf{Amplification cycles persist} — Synthetic signals deform
    identity faster than it can be rebuilt.

    \item \textbf{Ecological displacement is irreversible} — Synthetic species
    occupy identity regions.

    \item \textbf{Boundary anchors cannot be restored} — Identity lacks structural
    enclosure.

    \item \textbf{Restorative deficit exceeds threshold} — Identity curvature
    cannot regenerate fast enough (G6.6).
\end{itemize}

\medskip

\noindent When Identity Re-Anchoring fails, restorative deficit becomes
structural. Shared floors destabilize (G6.2), boundaries weaken (G6.3), deficit
magnitude increases (G6.5), and restoration thresholds (G6.6) are crossed.
Identity Re-Anchoring is the geometric indicator that identity curvature is being
restored—and the first signal when restoration becomes insufficient.

\subsection*{G6.5 Restorative Deficit Quantification}

\noindent Restorative Deficit Quantification formalizes the measurement of how
far human-domain restorative processes fall behind synthetic deformation. A
restorative deficit is geometric: it is the difference between the curvature
required to restore stability and the curvature actually regenerated by
human-domain systems. Quantification provides the manifold-level metric needed to
predict restoration failure, collapse thresholds, and recovery feasibility.

\medskip

\noindent In stable conditions, restorative capacity exceeds deformation load.
Rhythms re-synchronize (G6.1), shared floors reconstruct (G6.2), boundaries
reinforce (G6.3), and identity re-anchors (G6.4). A deficit does not emerge
unless synthetic pressure fields (G3), dominance modes (G5), and institutional
deformation modes (G4) outpace restorative geometry.

\medskip

\noindent Restorative Deficit Quantification becomes necessary when:

\begin{itemize}[leftmargin=1.2cm]
    \item \textbf{Restoration lags deformation} — Synthetic pressure increases
    faster than restoration can regenerate curvature.

    \item \textbf{Multiple deformation modes compound} — Boundary failure,
    role collapse, rhythm fragmentation, meaning leakage, and container failure
    co-occur (G4.2–G4.6).

    \item \textbf{Dominance modes intensify} — Speed, compliance, incentive,
    amplification, displacement, or dominant signals exceed ecological tolerances
    (G5.1–G5.6).

    \item \textbf{Recovery bands collapse} — Slack reserves fall below critical
    thresholds (G2.6).

    \item \textbf{Identity drift accelerates} — Identity curvature becomes
    unstable (G6.4).
\end{itemize}

\medskip

\noindent Restorative Deficit Quantification measures the deficit through five
geometric indicators:

\begin{itemize}[leftmargin=1.2cm]
    \item \textbf{Curvature Deficit} — The difference between required curvature
    and regenerated curvature.

    \item \textbf{Rhythmic Deficit} — The gap between temporal deformation and
    re-synchronization capacity.

    \item \textbf{Boundary Deficit} — The shortfall in boundary curvature needed
    to re-establish identity enclosures.

    \item \textbf{Identity Deficit} — The displacement between identity drift
    vectors and identity re-anchoring capacity.

    \item \textbf{Container Deficit} — The gap between container failure modes
    and containment restoration capacity.
\end{itemize}

\medskip

\noindent Restorative Deficit Quantification deforms the meaning manifold by:

\begin{itemize}[leftmargin=1.2cm]
    \item \textbf{Revealing collapse trajectories} — Deficit magnitude predicts
    which collapse paths (G7) will activate.

    \item \textbf{Reducing restoration feasibility} — Large deficits indicate
    restoration cannot keep pace with deformation.

    \item \textbf{Accelerating drift} — Deficit growth increases drift vector
    magnitude (G2.5).

    \item \textbf{Weakening ecological resilience} — Deficits reduce the
    manifold’s ability to absorb synthetic pressure.

    \item \textbf{Destabilizing identity} — Identity curvature becomes
    increasingly difficult to restore.
\end{itemize}

\medskip

\noindent Restorative Deficit Quantification is amplified by:

\begin{itemize}[leftmargin=1.2cm]
    \item \textbf{Synthetic velocity dominance} — High-speed deformation outruns
    restoration (G5.1).

    \item \textbf{Amplification curvature} — Algorithmic amplification increases
    deformation magnitude (G5.4).

    \item \textbf{Ecological displacement} — Synthetic species occupy restorative
    niches (G5.5).

    \item \textbf{Counterfeit invariants} — False stability masks deficit growth
    until thresholds are crossed.

    \item \textbf{Cross-substrate coupling} — Deficits propagate across human,
    digital, machine, and institutional layers.
\end{itemize}

\medskip

\noindent Restorative Deficit Quantification fails when:

\begin{itemize}[leftmargin=1.2cm]
    \item \textbf{Deficit exceeds critical threshold} — Restoration becomes
    geometrically impossible (G6.6).

    \item \textbf{Synthetic curvature becomes dominant} — Restoration cannot
    regenerate curvature against dominant signals (G5.6).

    \item \textbf{Recovery bands collapse} — Slack reserves fall below
    regenerative minimums (G2.6).

    \item \textbf{Identity anchors dissolve} — Identity curvature cannot be
    re-established.

    \item \textbf{Container failure becomes systemic} — Containment cannot be
    restored (G4.6).
\end{itemize}

\medskip

\noindent When Restorative Deficit Quantification fails, restoration becomes
non-viable. Thresholds are crossed (G6.6), collapse paths activate (G7), and the
manifold enters a regime where synthetic pressure governs structural dynamics.
Restorative Deficit Quantification is the geometric indicator that restoration is
falling behind—and the final signal before restoration becomes impossible.

\subsection*{G6.6 Restoration Failure Thresholds}

\noindent Restoration Failure Thresholds formalize the geometric limits beyond
which human-domain restoration cannot regenerate curvature fast enough to
counteract synthetic deformation. A threshold is not a metaphor or heuristic; it
is a structural boundary in the manifold where restorative capacity falls below
the minimum curvature required for stability. Once crossed, restoration becomes
non-viable, and collapse paths (G7) activate.

\medskip

\noindent In stable conditions, restorative processes exceed deformation load.
Rhythms re-synchronize (G6.1), shared floors reconstruct (G6.2), boundaries
reinforce (G6.3), identity re-anchors (G6.4), and restorative deficit remains
below critical magnitude (G6.5). Failure thresholds do not appear unless
synthetic pressure fields (G3), dominance modes (G5), and institutional
deformation modes (G4) compound.

\medskip

\noindent Restoration Failure Thresholds emerge when:

\begin{itemize}[leftmargin=1.2cm]
    \item \textbf{Curvature deficit exceeds critical magnitude} — Required
    curvature surpasses regenerative capacity (G6.5).

    \item \textbf{Temporal deficit becomes irreversible} — Rhythmic
    re-synchronization cannot restore continuity bands (G6.1).

    \item \textbf{Shared floors cannot be rebuilt} — Interpretive surfaces fail
    to regain curvature (G6.2).

    \item \textbf{Boundary curvature collapses} — Identity boundaries cannot be
    re-established (G6.3).

    \item \textbf{Identity anchors dissolve} — Identity curvature cannot be
    restored (G6.4).

    \item \textbf{Container failure becomes systemic} — Containment architecture
    cannot be regenerated (G4.6).
\end{itemize}

\medskip

\noindent Restoration Failure Thresholds deform the meaning manifold by:

\begin{itemize}[leftmargin=1.2cm]
    \item \textbf{Eliminating restorative feasibility} — Restoration becomes
    geometrically impossible.

    \item \textbf{Accelerating drift} — Drift vectors increase in magnitude and
    directionality (G2.5).

    \item \textbf{Collapsing recovery bands} — Slack reserves fall below
    regenerative minimums (G2.6).

    \item \textbf{Destabilizing identity} — Identity curvature becomes
    uncontainable.

    \item \textbf{Opening collapse paths} — Threshold crossing activates collapse
    geometries formalized in G7.
\end{itemize}

\medskip

\noindent Restoration Failure Thresholds are amplified by:

\begin{itemize}[leftmargin=1.2cm]
    \item \textbf{Synthetic velocity dominance} — High-speed deformation outruns
    all restorative modes (G5.1).

    \item \textbf{Amplification curvature} — Algorithmic amplification increases
    deformation magnitude (G5.4).

    \item \textbf{Ecological displacement} — Synthetic species occupy restorative
    niches (G5.5).

    \item \textbf{Dominant signals} — Synthetic curvature becomes ecological law
    (G5.6).

    \item \textbf{Counterfeit invariants} — False stability masks threshold
    approach until collapse is underway.
\end{itemize}

\medskip

\noindent Restoration Failure Thresholds are crossed when:

\begin{itemize}[leftmargin=1.2cm]
    \item \textbf{Restorative deficit becomes supercritical} — Deficit magnitude
    exceeds regenerative capacity.

    \item \textbf{Curvature cannot be regenerated} — Restoration cannot rebuild
    curvature faster than synthetic pressure deforms it.

    \item \textbf{Identity cannot be re-anchored} — Identity curvature loses all
    structural anchors.

    \item \textbf{Containers cannot be restored} — Containment architecture
    collapses beyond repair.

    \item \textbf{Synthetic curvature governs the manifold} — Dominant signals
    enforce machine-native geometry.
\end{itemize}

\medskip

\noindent When Restoration Failure Thresholds are crossed, the manifold enters a
collapse regime. Restoration becomes non-viable, synthetic pressure governs
structural dynamics, and collapse paths (G7) activate. Restoration Failure
Thresholds are the geometric indicator that the restorative stack has reached its
limit—and that the manifold is transitioning into collapse geometry.

\medskip

\noindent G6.6 closes the Restorative Deficit Dynamics layer and provides the
structural foundation for the collapse geometries formalized in G7.
